\chapter{Forskinga utan objekt}

\begin{motto}
Eit fag som har gjort si eiga manglande etterprøvbarheit til ein metode,\\ har slutta å vere eit fag som kan ta feil.
\end{motto}

Då designfaget skulle inn i universitetet for fullt — i Storbritannia frå 1960-talet, breiare frå 1990-talet — møtte det eit krav det ikkje kunne oppfylle: kravet om forsking. Universitetet er bygt kring frambringing av etterprøvbar kunnskap. Faget hadde to ærlege val: byggje eit fundament som gjorde forsking mogleg, eller innrømme at det er eit handverksfag og be om ein annan plass. Det gjorde verken det eine eller andre. Det fann ein tredje veg — å omdefinere forsking slik at det det alt gjorde, kunne telje som forsking. Resultatet går under namnet \emph{forsking gjennom design}.

\section{Frayling og dei tre preposisjonane}

Den kanoniske inndelinga kom frå Christopher Frayling i 1993 \autocite{frayling1993}. Han skilde mellom forsking \emph{om} design (designhistorie og -teori), forsking \emph{for} design (kunnskap framskaffa til bruk i designarbeid), og forsking \emph{gjennom} design (kunnskap produsert i sjølve handlinga å designe). Dei to fyrste er uproblematiske, men er som regel historie, samfunnsvitskap eller materialvitskap brukt på design — ikkje ein eigen designvitskap. Den tredje var den faget hengde framtida si på, fordi berre den gjer sjølve designpraksisen til forsking. Seinare arbeid har raffinert inndelinga \autocite{frayling1993b}, men grunnstrukturen står: artefaktet pluss refleksjonen kring det skal utgjere eit kunnskapsbidrag.

For at dette skal vere forsking, må artefaktet eller prosessen produsere ein påstand om verda som andre kan ta opp, etterprøve, byggje på, eller vise er feil. Det er minstekravet — ikkje at design skal vere som fysikk, berre at det skal vere kumulativt og etterprøvbart, slik at feltet lærer noko som varer ut over den einskilde forskaren. William Gaver, sjølv ein framståande forskar i feltet, stilte spørsmålet ærleg: kva bør vi eigentleg vente oss av forsking gjennom design? \autocite{gaver2012what} Svaret hans — at feltet produserer «mellombels, provisoriske» kunnskapsformer snarare enn teoriar — er ærleg, men det er òg ei innrømming. Bruce Archer hadde tjue år tidlegare hevda at design var ein eigen disiplin med eit eige kunnskapsgrunnlag \autocite{archer1979discipline}; Gavers forsiktige formulering måler avstanden mellom den ambisjonen og det feltet faktisk leverte.

\section{Kvifor det ikkje akkumulerer}

Sjå på kva eit typisk arbeid leverer. Ein forskar designar eit artefakt — eit produkt, ein teneste, eit spekulativt objekt — og skriv ein tekst om kva han lærte. Spør: kva er påstanden? Kva ville telje som å motbevise han? Kan ein annan etterprøve han? I det overveldande fleirtalet av tilfelle finst ingen avgrensa påstand, ingenting ville telje som motbevis, og ingen kan etterprøve noko, fordi resultatet er bunde til den einskilde forskaren sin praksis, kontekst og person. Arbeidet er ofte rikt og gjennomtenkt. Men det er ikkje kumulativt. Den neste forskaren byrjar ikkje der den førre slutta; ho byrjar på nytt, med sitt eige artefakt, sin eigen refleksjon. Etter tiår med forsking gjennom design finst det ingen veksande kropp av etablerte resultat som studentane lærer som grunnlag.

Og her lukkar sirkelen seg mot førre kapittel. I staden for å sjå denne mangelen som ein veikskap, har faget gjort han til ein dyd ved hjelp av det lånte omgrepet om vrange problem \autocite{buchanan1992wicked}. Forskinga treng ikkje vere kumulativ, fordi problema er vrange; resultata treng ikkje vere etterprøvbare, fordi kvart tilfelle er unikt. Schöns reflekterande utøvar \autocite{schon1983} blir løfta frå pedagogikk til epistemologi: refleksjonen \emph{er} kunnskapen. Men ein refleksjon som ingen annan kan etterprøve, er ikkje kunnskap i den forstand ein doktorgrad er meint å sertifisere. Han er ei velformulert erfaring.

\section{Doktorgraden som ritual}

Konsekvensen ser ein tydelegast i den praksisbaserte doktorgraden, som vaks fram frå 1990-talet. Ein kandidat brukar år på eit kreativt arbeid, supplerer det med ein refleksiv tekst, og får tittelen doktor. Arbeida er ofte framifrå \emph{som arbeid}. Men ein doktorgrad er meint å sertifisere at innehavaren har produsert eit etterprøvbart, originalt bidrag til ein kumulativ kunnskapskropp. Når kunnskapskroppen ikkje er kumulativ og bidraget ikkje er etterprøvbart, sertifiserer tittelen ingenting av dette. Han sertifiserer at kandidaten gjorde eit seriøst kreativt arbeid og reflekterte godt over det — ein heider verdt å gje, men ein annan heider enn den tittelen lovar. Faget brukar den gamle tittelen for den nye tingen fordi den gamle ber den autoriteten faget vil ha. Det er den same dobbeltbokføringa som går att gjennom heile boka: faget krev forskaren sin status for handverkaren sitt arbeid.

Dette er ikkje eit åtak på den einskilde forskaren, som ofte arbeider med stor integritet innanfor reglane feltet har gjeve. Det er ein diagnose av reglane. Eit fag som hadde eit fundament, kunne late forsking gjennom design bli ekte forsking: artefaktet kunne berast av falsifiserbare påstandar om kva forma gjer, kvifor ho tek den forma ho tek, kva ho ber og kva ho legg til side. Det er nett det fundamentet faget manglar, og difor er forsking gjennom design dømd til å sirkle kring si eiga uangripelegheit.

\section{Den britiske maskinen: korleis ein administrativ reform skapte ein doktorgrad}

Det er lett å lese forsking gjennom design som ein idéhistorisk feilslutning — ein tankefeil nokon gjorde og andre arva. Det er ikkje heile sanninga, og det er ikkje den mest avslørande. For den praksisbaserte doktorgraden vart ikkje fyrst og fremst tenkt fram; ho vart \emph{administrert} fram, av ein bestemt institusjonell maskin i Storbritannia, og kjenner ein maskina, ser ein at heile konstruksjonen er driven av insentiv heller enn av epistemologi. Historia byrjar med ein lov. I 1992 omgjorde \textit{Further and Higher Education Act} dei britiske polyteknika — der det meste av kunst- og designutdanninga budde — til fullverdige universitet. Over natta måtte institusjonar som i generasjonar hadde drive handverksopplæring i atelier, levere det universitet leverer: forsking, doktorgrader, publikasjonar som kunne teljast. Den gamle eksamensstyresmakta for polyteknika, CNAA, hadde alt frå åttitalet opna for forskingsgrader i kunst og design; no vart den moglegheita til eit krav, fordi finansieringa hang i det.

Og finansieringa hang i ein bestemt mekanisme: den nasjonale forskingsevalueringa, fyrst kalla Research Assessment Exercise, seinare Research Excellence Framework. Med jamne mellomrom vert kvar institusjon vurdert, og pengar fordelt etter kor mykje «forsking» av høg kvalitet ho produserer. Eit designinstitutt som ikkje kunne vise fram forsking, ville misse pengar til eit som kunne. Dermed oppstod eit massivt, reint økonomisk press for å \emph{omklassifisere} det designarar alt gjorde — utstillingar, prototypar, porteføljar — som forskingsoutput. Christopher Frayling, som gav oss dei tre preposisjonane \autocite{frayling1993}, var sjølv rektor ved Royal College of Art og djupt inne i denne maskinen; inndelinga hans var ikkje berre ein filosofisk distinksjon, men eit verktøy for å gjere kunstpraksis teljande i eit system bygd for vitskap. Det same presset fødde «alternativ-format»-avhandlinga: ein doktorgrad der hovudbidraget er eit kreativt arbeid, supplert med ein skriftleg «exegesis» eller refleksjon. Sjå rekkjefølgja nøye. Behovet kom ikkje frå ei innsikt om at design produserer ein eigen kunnskap. Det kom frå eit krav om å telje, og teorien om forsking gjennom design vart bygd \emph{etterpå}, for å rettferdiggjere det tellinga alt kravde. Epistemologien er ein rasjonalisering av eit budsjett.

\begin{kasus}{REF og den motsette rekkjefølgja}
I eit sunt fag kjem rekkjefølgja slik: ein har eit kunnskapsobjekt, så byggjer ein ein doktorgrad for å sertifisere bidrag til det, så vurderer ein institusjonane etter kor godt dei bidreg. I den britiske designakademia gjekk det baklengs. Fyrst kom evalueringa som delte ut pengar etter «forsking» (RAE frå 1986, REF frå 2014). Så kom kravet om at designinstitutta måtte produsere noko som kunne kallast forsking for å overleve evalueringa. Så, til sist, kom teorien — Fraylings preposisjonar, exegesis-avhandlinga, «practice as research» — som forklarte korleis det designarar alt gjorde, kunne reknast som forsking. Objektet vart ikkje oppdaga; det vart postulert, fordi finansieringsmaskina kravde at eitt fanst. Dette forklarer ein observasjon som elles er gåtefull: kvifor forsking gjennom design er full av metarefleksjon om kva ho \emph{er} og kva ein bør \emph{vente} av henne \autocite{gaver2012what}, men fattig på etablerte resultat. Eit felt som var bygd kring eit objekt, ville bruke tida på objektet. Eit felt som er bygd for å tilfredsstille ei teljemaskin, brukar tida på å forsvare at det som vert talt, faktisk tel. Mengda av definisjonsarbeid er ikkje eit teikn på ungdom. Det er fingeravtrykket til ein disiplin som vart kravd fram før han hadde noko å forske på.
\end{kasus}

\section{Den nordiske kuren: å døype om mangelen til ein mellomkategori}

Når objektet manglar, men kravet om kunnskap står, oppstår ein føreseieleg sjanger: framlegg som gjev mangelen eit verdig namn. Den nordiske designforskinga, som tok dette på største alvor, har levert dei mest gjennomtenkte døma, og nett difor er dei mest opplysande. Jonas Löwgren føreslo omgrepet \emph{intermediate-level knowledge}: ein kunnskap som ligg mellom det einskilde artefaktet og den allmenne teorien, formidla mellom anna gjennom «annoterte porteføljar» — samlingar av designdøme med kommentarar som peikar på det generelle i det særskilde \autocite{lowgren2013intermediate}. Tanken er ærleg og nyttig: ein vedgår at forsking gjennom design ikkje produserer teori, og prøver å namngje kva ho då produserer. Men sjå kva namnet faktisk gjer. Det «mellomnivået» er definert ved kva det \emph{ikkje} er — ikkje eit reint artefakt, ikkje ein etterprøvbar teori — og det manglar nett det som ville gjere det til kumulativ kunnskap: eit kriterium for når éin annotert portefølje er rettare enn ein annan. To forskarar kan annotere same portefølje motstridande, og ingenting avgjer kven som har rett. Namnet er presist; tomrommet det namngjev, er det same.

Den same ærlege famlinga finn ein i den nordiske debatten om kva «doctorateness» eigentleg er i designfaga. Fredrik Nilsson og Halina Dunin-Woyseth ved Arkitektur- og designhøgskolen i Oslo har skrive direkte om problemet: korleis sertifisere ein doktorgrad i eit praksisfelt der dei vanlege kriteria for vitskapleg bidrag ikkje passar? Svaret deira er å forhandle fram noko mellom «kjennarskap» (connoisseurship) og «kritikk» \autocite{nilsson2012doctorateness} — å lene seg på den røynde utøvaren si vurdering av kvalitet, men kle henne i akademiske ord. Dunin-Woyseth innførte òg, frå vitskapssosiologien, skiljet mellom «mode 1» og «mode 2»-kunnskap for å gje praksisforskinga ein plass \autocite{duninwoyseth2011making}. Dette er det beste faget har; det er reflektert, sjølvkritisk og institusjonelt seriøst. Og likevel endar det på same staden som alt anna i boka: kjennarskap er det trente blikket frå atelieret, no løfta til doktorgradsnivå, men framleis utan ein uavhengig målestokk under seg. Joep Frens si avhandling frå Eindhoven, der den rike samhandlinga mellom form, funksjon og interaksjon vert utforska gjennom sjølve det å designe \autocite{frens2006rich}, er eit framifrå døme på sjangeren — og samstundes eit døme på at sjølv dei beste avhandlingane sertifiserer eit dugande kreativt arbeid med refleksjon, ikkje eit etterprøvbart bidrag til ein veksande kunnskapskropp.

Trekk no saman det dette kapitlet har vist. Den britiske maskina forklarer kvifor doktorgraden vart kravd fram — av eit budsjett, ikkje av eit objekt. Den nordiske kuren forklarer kva faget gjer når objektet likevel manglar — det døyper mangelen om til ein verdig mellomkategori og lener seg til slutt på kjennarens blikk. Begge er ærlege på sin måte; ingen av dei skaffar det som manglar. Og det er heile poenget. Forsking gjennom design er ikkje eit dårleg utført forsøk på ekte forsking som betre folk kunne ha fiksa. Det er eit institusjonelt apparat bygd presist for å produsere forskingens form — grader, tidsskrift, evalueringspoeng — utan forskingens substans, fordi substansen krev eit objekt faget ikkje har. Maskina går utmerkt. Ho lagar berre ikkje kunnskap som hopar seg opp.

\section{Disputasen: artefaktet på veggen}

Lat oss setje oss i salen, for ein praksisbasert disputas er ein scene som fortel meir enn nokon definisjon. Det er ein ettermiddag i eit auditorium ved ein britisk kunst- og designhøgskule, ein gong på 2000-talet. På veggane heng kandidaten sitt arbeid: ein serie objekt, kanskje møblar, kanskje interaktive prototypar, kanskje spekulative ting som ikkje skal seljast men tenkjast med. På bordet ligg ein tekst, «the exegesis», nokre titals tusen ord der kandidaten reflekterer over kva ho gjorde og kva ho lærte. Kommisjonen sit framme — to, kanskje tre granskarar, ein intern og ein ekstern, kledde i si vante akademiske alvor. Dei har lese teksten. Dei har sett på objekta. Og no skal dei gjere det ein kommisjon alltid gjer: finne fram til \emph{påstanden}, det avgrensa, originale bidraget til kunnskap som ein doktorgrad skal sertifisere.

Sjå kva som skjer i denne augneblinken, for det er her heile kapitlet kjem til syne i éin gest. Granskaren lener seg fram og spør, høfleg og uunngåeleg: «Kva er det nye her, som andre kan byggje vidare på?» Kandidaten svarar — ærleg, gjennomtenkt — med å skildre prosessen sin, kva ho oppdaga undervegs, korleis materialet motsette seg henne, korleis ein ny måte å sjå på vaks fram. Det er eit godt svar om eit godt arbeid. Men det er ikkje svaret på spørsmålet. For spørsmålet galdt påstanden: kva seier dette arbeidet om verda som ein annan forskar kan ta opp, etterprøve, eller vise er feil? Og i det overveldande fleirtalet av slike disputasar finst det ikkje noko slikt. Det finst objekt, det finst innsikt, det finst ein person som har lært mykje. Kommisjonen leitar etter ein påstand på veggen og finn ein portefølje. Dei veit det godt, dei fleste av dei; difor flyttar samtalen seg, nesten umerkeleg, frå «er dette sant» til «er dette godt utført» — frå etterprøving til kjennarskap\autocite{nilsson2012doctorateness}. Og når dei to timane er omme, gjev dei kandidaten tittelen doktor. Dei sertifiserer eit framifrå kreativt arbeid med refleksjon, slik Joep Frens si Eindhoven-avhandling om den rike samhandlinga mellom form og interaksjon er eit framifrå døme på\autocite{frens2006rich}. Dei sertifiserer ikkje eit etterprøvbart bidrag til ein veksande kunnskapskropp, fordi det ikkje var noko slikt å sertifisere. Ritualet er fullendt. Forma er forskingens; substansen er handverkets.

\section{Frayling ved skrivebordet, 1993}

Spol attende til kjelda. Det er 1993, og ved Royal College of Art i London sit rektoren, Christopher Frayling — kunsthistorikar, fjernsynsmann, mannen som seinare skulle bli adla — og skriv eit kort arbeidsnotat for skulens eigen serie, \textit{Royal College of Art Research Papers}. Tittelen er beskjeden: «Research in Art and Design»\autocite{frayling1993}. Notatet er ikkje langt, og det opnar ikkje med epistemologi, men med ein scene frå filmen \textit{Frankenstein} — galne vitskapsmannen i laboratoriet — som eit bilete på kva folk \emph{trur} forsking er. Frayling vil rydde unna karikaturen og vise at også kunstnaren og designaren forskar. Og så, nesten i forbifarten, kjem distinksjonen som skulle forme eit heilt felt i tiåra etter: forsking \emph{om}, \emph{for} og \emph{gjennom} kunst og design.

Det er verdt å hugse kva slag tekst dette eigentleg er. Det er ikkje eit filosofisk verk, ikkje resultatet av tjue års grunnlagsarbeid; det er eit framifrå formulert posisjonsnotat skrive av ein rektor i ein institusjon som nett då var i ferd med å bli pressa inn i den nasjonale forskingsevalueringa. Frayling sjølv var med på å reformere kunst- og designforskinga i Storbritannia; han visste betre enn dei fleste at polyteknika nyleg var gjorde til universitet, at pengar no hang i evne til å vise fram «forsking», og at nokon måtte gje eit namn til det kunstnarar og designarar alt gjorde, slik at det kunne teljast. Dei tre preposisjonane var det namnet. Det var ein elegant tankegang, og han har tent feltet på utallige vis. Men sjå rekkjefølgja: distinksjonen vart ikkje destillert ut av eit modent kunnskapsobjekt; han vart formulert for å gje eit institusjonelt behov ei verdig form. Det at den mest siterte definisjonen av designforsking opphavleg var eit fjorten-siders notat med ein \textit{Frankenstein}-vits i opninga, er ikkje eit teikn på at feltet var ungt og lovande. Det er fingeravtrykket til eit objekt som vart postulert, ikkje oppdaga.

\begin{kasus}{Panelet som skåra ei utstilling}
Tenk på rommet der avgjerda faktisk fall — ikkje disputasen, men evalueringspanelet. Under Research Assessment Exercise, og seinare Research Excellence Framework, sat ekspertpanel og skåra forskingsoutput frå kvar institusjon på ein skala, og pengane følgde poenga. For dei fleste fag var outputen ein artikkel: panelet las påstanden, vurderte metoden, sjekka om funnet heldt. Men for kunst og design måtte panelet skåre noko anna — ei utstilling, ein katalog, ein serie objekt, ein performance — som «forsking». Sjå kva det krev. Eit panel som skal gje ein togradig dom over om ei utstilling representerer «verdsleiande» eller berre «internasjonalt framifrå» forsking, må ha eit kriterium for kva som skil forsking frå framifrå kunst. Det kriteriet finst ikkje. Så panelet gjorde det einaste det kunne: det las den vedlagde «300-ords-skildringa» som kvar output måtte ha med seg — eit lite stykke prosa der innsendaren forklarte kva forskingsspørsmål, kva metode og kva «insight» arbeidet representerte. Med andre ord: kunstverket åleine kunne ikkje skårast som forsking; det måtte følgjast av ein tekst som \emph{omtala} det som forsking, og det var teksten panelet eigentleg vurderte. Her ser ein heile konstruksjonen i miniatyr. Objektet ber ikkje påstanden; ein påklistra refleksjon ber han\autocite{gaver2012what}. Pengane vart fordelte, institusjonane overlevde, og kvar runde produserte fleire hundre slike skildringar — eit veksande arkiv av tekstar som forklarer at det som vart talt, faktisk tel. Det var akkurat så mykje forsking som maskina kravde, og ikkje eit gram meir.
\end{kasus}

\vignett

Vi har no følgt faget gjennom metoden som kollapsa, vitskapsdraumen som ikkje feste seg, det vrange som vart alibi, atelieret som overfører smak, og forskinga som ikkje akkumulerer. Med Del \textsc{v} forlèt vi den akademiske ambisjonen og følgjer faget dit det faktisk lukkast med å spreie seg: ut av seg sjølv, inn i meininga, inn i skjermen, og inn i økonomien.

\laeringsmal{\begin{itemize}
\item Gjere greie for Fraylings tredeling (forsking om, for og gjennom design) og kvifor «gjennom» er den problematiske kategorien.
\item Forklare korleis REF/RAE-insentiva i britisk høgre utdanning skapte den praksisbaserte doktorgraden — epistemologien som rasjonalisering av eit budsjett.
\item Vurdere kva ein doktorgrad er meint å sertifisere, og kvifor forsking gjennom design ikkje akkumulerer.
\end{itemize}}

\ovingar{\begin{enumerate}
\item Finn ein praksisbasert ph.d. (artefakt pluss refleksjon). Kva er den avgrensa, etterprøvbare påstanden? Kva ville telje som å motbevise han?
\item Gaver kallar kunnskapen frå forsking gjennom design «provisorisk». Er det ei ærleg innrømming eller eit forsvar? Drøft.
\item Kva ville krevjast for at forsking gjennom design skulle bli kumulativ? Skisser eit konkret, fleirårig forskingsprogram som ville oppfylle kravet.
\end{enumerate}}

\vidarelesing{\fullcite{frayling1993}; \fullcite{gaver2012what}; \fullcite{archer1979discipline}; \fullcite{schon1983}; \fullcite{cross2006}.}
